\section{Scaling out and connecting to data}
\label{sec:scale}

\modelly{} is not limited to one process on one CPU, and connects a model to arriving data, to
continuous fields and to the browser. For each capability this section
gives the mechanism, what was measured and where it falls short; \cref{tab:scale}
collects the numbers, and \cref{app:scale} has the designs, the measurement conditions
and the file behind every one of them. The measurements were made on the laptop of
\cref{sec:performance} while it was shared with other work and without the fixed frequency of
\cref{sec:playbooks}, and \cref{sec:threats} says which are repeated runs and which are single
runs; the conclusions that matter (byte-identical outputs, order-of-magnitude speed-ups,
shortfalls by a factor) do not turn on the last digit of a timing.

\begin{table}[tbp]
  \centering
  \footnotesize
  \caption{Scaling out and connecting to data: what was measured and where it falls short.
  Details and sources in \cref{app:scale}.}
  \label{tab:scale}
  \begin{tabularx}{\linewidth}{>{\raggedright\arraybackslash}p{0.17\linewidth}>{\raggedright\arraybackslash}X>{\raggedright\arraybackslash}X}
    \toprule
    Capability & Measured & Falls short \\
    \midrule
    Split runs & Byte-identical to one process at 1, 2, 4 and 16 partitions on three models;
    MPI with 2 and 4 ranks; Dask with 2 workers & Weak scaling 72, 32 and 17\% at 2, 4 and 8
    workers with 10,000 people per worker, 73, 46 and 9\% with 40,000 (single runs on a loaded
    laptop); runs at 12 and 16 workers swamped by contention; 64 workers unmeasured; Ray, and
    MPI on macOS and Windows, untested \\
    \addlinespace
    GPU backend & Steps 34.9 times faster at 10 million agents, identical integer state &
    Construction not faster (1.3 to 2.3 times in total); slower at 20,000 agents, the
    crossover unmeasured; cohort tables 200 times slower; Metal untested \\ \addlinespace
    Differentiable models & Two zoo models' parameters recovered within 2\% and 7\%; pathwise
    gradients within $10^{-6}$ and $10^{-4}$ of finite differences & Gradients through discrete
    choices are approximations; the straight-through estimator is biased \\ \addlinespace
    Data assimilation & Nowcast error 5.07 against 69.2 unfiltered; contact rate 0.451 for
    0.45; bit-identical for a seed & One dataset and one seed; brokers tested with fake
    clients; replicates step in one process \\ \addlinespace
    Fields and coupling & Mass conserved to $10^{-12}$; heat kernel to about $10^{-6}$; two
    zoo models with closed budgets & Orthogonal grids with uniform spacing only; upwind
    advection smears fronts \\ \addlinespace
    Emulators & Parameter box reduced from 17.1 to 2.5\% in four waves, truth inside & Half the
    runs of a hypercube only on the zoo model (factors 1.3 to 2.7); Branin 1.1 to 1.4, Ishigami
    0.8 to 1.4 \\ \addlinespace
    Synthetic populations & Totals and ages exact, sexes within 1\%; the SEIR model's final
    size within 0.05; three real-data zoo models & Made-up marginals; no real joint
    distribution to test against \\ \addlinespace
    Browser & 62 of 62 runs: 31 models, two wheels & First load 8.0, 12.3 and 21.4~s by line
    and model, under 10~s only without \code{scipy} at 50~Mbit/s \\
    \bottomrule
  \end{tabularx}
\end{table}

\subsection{Per-agent random streams and split runs}
\label{sec:split}

A model can give every agent random generators of its own, seeded from the scenario's seed
and the agent's key, so that what an agent draws no longer depends on the activation order or
on the population (the flag is off by default). One run can then be split across local, MPI,
Dask or Ray workers (a grid into blocks with halos, a network by breadth-first banding, a
continuous space by recursive bisection), with ghost agents exchanged before a phase that
reads across a border and movers migrating after a phase that moves them. On three reference
models every reporter and every agent's state is byte-identical to the single-process run at
1, 2, 4 and 16 partitions, under MPI with 2 and 4 ranks and on two Dask workers
(\cref{app:split}). The weak scaling of \cref{tab:scale} comes from single runs on a laptop
shared with other jobs, whose hybrid processor is the main limit: eight copies of the same
work with no communication at all already run at 24\% of one, because beyond the four
performance cores the copies land on efficiency cores, and the protocol keeps 51 to 82\% of
what the machine gives up to eight workers; at 12 and 16 workers the split runs took about a
hundred times longer per step, under contention a single run cannot separate from the
protocol's cost, so no efficiency is claimed there (\cref{tab:scaling}). The aim of 70\% efficiency at 64
workers could not be measured here; the Ray backend, and MPI on macOS and Windows, are
untested.

\subsection{A GPU backend}
\label{sec:gpu}

Column phases, cohort tables and the cell-space operators run on CuPy or JAX with a
\code{numpy} fallback, from counter-based (Philox) random streams that draw the same numbers
on every backend. At ten million agents the grid epidemic's 60 steps take 0.67~s on an
RTX~5050 laptop GPU against 23~s for the CPU column variant, 34.9 times faster (median of three
rounds, 31.9 to 50.9), with the same infected series and the same digest of the final state on \code{numpy}
and on CuPy (\cref{app:gpu}). Building ten million agent objects takes 28 to 48~s on either
side, so build plus steps gains only 1.3 to 2.3 times; at 20,000 agents the GPU is slower
than the CPU (0.16 to 0.24~s against 0.08 to 0.13~s) and the crossover was not measured; a cohort
table of under 200 rows is about 200 times slower, because kernel launches cost more than its
arithmetic; and Apple Metal was not measured.

\subsection{Differentiable models}
\label{sec:diff}

Column phases traced by JAX give gradients of a model's outputs with respect to its
parameters, exact through normal and uniform draws and approximate through Bernoulli choices
(straight-through, Gumbel-softmax or score-function estimators, each warning that it is one).
Gradient calibration recovers two zoo models' parameters from runs of the agent-based
originals, and pathwise gradients agree with finite differences (\cref{tab:scale},
\cref{app:diff}); the straight-through estimator is biased and fails that check, by design, so
a gradient through a discrete choice is an approximation whose use the documentation delimits.

\subsection{Live data and data assimilation}
\label{sec:assimilation}

Observations arrive from a replayed table, a growing CSV file, an HTTP endpoint or a message
broker; an ensemble of replicates advances in lock step, a bootstrap particle filter
\citep{gordon1993novel} weights and resamples it, and an ensemble Kalman filter
\citep{evensen2003ensemble} adjusts continuous parameters, both bit-identical for a seed. On
the tutorial's synthetic town the filtered nowcast of the unobserved number infected has a
root-mean-square error of 5.07 people against 69.2 for the unfiltered ensemble, and the
Kalman filter recovers the contact rate as 0.451 for a truth of 0.45 (\cref{app:assimilation});
that is one dataset and one seed, not a coverage study. Message brokers were tested against
fake clients only, and the replicates step in one process.

\subsection{Fields and coupling}
\label{sec:couple}

A field is a property layer with a solver (diffusion, upwind advection and decay, in explicit
sub-steps that refuse an unstable time step, or backward Euler), coupled to differential
equations and stock-flow components under a schedule that is checked when the model is built.
Closed and periodic boundaries conserve mass to a relative $10^{-12}$ and a spreading Gaussian
matches the heat kernel to about $10^{-6}$; two zoo models, foragers on a diffusing nutrient
and an epidemic coupled to hospital capacity, close their budgets to round-off and match their
analytical checks (\cref{app:couple}). Fields need orthogonal grids with uniform spacing, and
upwind advection smears sharp fronts.

\subsection{Emulators, active learning and history matching}
\label{sec:emulate}

A Gaussian-process or neural emulator trained on an experiment's runs, active learning that
picks the next runs by integrated variance reduction, and history matching in waves
\citep{vernon2010galaxy, andrianakis2015bayesian}. On the age-structured SEIR model four waves
of 20 runs reduce the parameter box to 2.5\%, with the true parameters inside after every wave
(\cref{app:emulate}). Active learning saves the runs it was meant to save on the zoo model and
not on the Branin and Ishigami functions, where a well-spread design is already close to the
best that a variance criterion blind to the outputs can do.

\subsection{Synthetic populations and models on real data}
\label{sec:population}

Iterative proportional fitting, integerisation, households, places within a distance limit and
a contact network with household, school and work layers build a starting population from
aggregate tables; no personal data is bundled, and three small real samples (OpenStreetMap
streets, ESA WorldCover land cover, census tracts) ship with their licences. The synthetic
population reproduces its marginals and starts the SEIR model within its final-size check, and
three zoo models run on the real samples against known answers (\cref{app:population}); the
marginals are made up, so no real joint distribution was tested against.

\subsection{\modelly{} in the browser}
\label{sec:browser}

\cmd{modelly export-web} writes a static site that runs a model with Pyodide in a browser tab,
with Studio's parameter panel, a space view and a reporter plot. The 14 bundled examples
and the 17 zoo models each ran from their export in headless Chromium with the
pure-Python wheel and with the compiled one, 62 of 62 runs (\cref{app:browser}). A cold first
visit is ready after 8.0~s on a 50~Mbit/s line for a model that needs no \code{scipy}, 12.3~s
for one that does, and 21.4~s on a Fast~4G profile, so the aim of under ten seconds holds only
on the faster line without \code{scipy}: Pyodide, \code{numpy} and \code{pandas} alone are
13.5~MB.
